\chapter{Approximate Counting in the Polynomial Hierarchy}
\label{ch:approx}

If exact counting is above the polynomial hierarchy, by Toda's theorem, approximate counting is within it. Stockmeyer\index{Stockmeyer's theorem} proved in 1983 that every $\sharpP$ function can be approximated to within a factor $(1 \pm \epsilon)$ by a randomized polynomial time algorithm with an $\NP$ oracle, and Sipser\index{Sipser Lautemann theorem}'s hashing technique, which underlies the proof, is a cousin of the vertical walk. This chapter presents those results and the fully polynomial schemes of Karp\index{Karp, Richard} and Luby for the special case of disjunctive normal form, and places them against the thesis.

\section{Stockmeyer's theorem}

\begin{theorem}[Stockmeyer]
Let $f \in \sharpP$. There is a randomized algorithm which, given $x$ and $\epsilon > 0$, with access to an $\NP$ oracle, runs in time polynomial in $|x|$ and $1/\epsilon$ and outputs a number within a factor $(1 \pm \epsilon)$ of $f(x)$ with probability at least $\tfrac34$. In particular, approximate counting is in $\BPP^{\NP}$, and hence in $\Sigma_3^p$.
\end{theorem}

The proof uses pairwise independent hashing. Let $S \subseteq \{0,1\}^m$ be the set to be counted, recognised by a polynomial time verifier. For a hash function $h : \{0,1\}^m \to \{0,1\}^k$ drawn from a pairwise independent family, such as the affine maps $h(r) = Ar + b$ over $\mathrm{GF}(2)$, the number of elements of $S$ with $h(r) = 0$ has mean $|S| 2^{-k}$ and variance at most the mean. If $|S| 2^{-k}$ is large, say at least $4/\epsilon^2$, then by Chebyshev the count $|S \cap h^{-1}(0)|$ is within a factor $(1 \pm \epsilon)$ of $|S| 2^{-k}$ with constant probability. The question whether $|S \cap h^{-1}(0)| \ge t$ is in $\NP$ for any polynomially bounded threshold $t$, since the witness is $t$ distinct elements of $S$ in the fibre, so the $\NP$ oracle can determine the count in the fibre by binary search. Trying $k = 0, 1, 2, \dots$ until the fibre count is in the right range, and multiplying by $2^k$, gives the estimate.

The structure is a vertical walk in disguise. The hash function slices $S$ into $2^k$ fibres of roughly equal size, which is the same as drawing a uniform point of $S$ and asking how many other points fall in its fibre. Each increment of $k$ halves the fibre, which is a ladder with ratio $\tfrac12$, and the number of rungs is $\log_2 |S| = O(m)$. The $\NP$ oracle replaces the level set sampler\index{level set sampler}: it answers the question whether a fibre has at least $t$ elements, which is a question about the width of a step, directly, by nondeterministic guessing, where the sampler would answer it by Monte Carlo. The theorem says that with an $\NP$ oracle every $\sharpP$ staircase can be climbed; the thesis says that without one the staircase of a self reducible problem can be climbed if and only if its level sets can be sampled.

\section{Sipser's hashing and the Valiant Vazirani lemma}

The hashing technique originates in Sipser's proof that $\BPP$ is in the polynomial hierarchy, in the version before Lautemann's simplification, and in the Valiant\index{Valiant, Leslie} Vazirani isolation lemma of Chapter~\ref{ch:pseudo}. In all three the idea is that a random linear map over $\mathrm{GF}(2)$ behaves, on any fixed set $S$, like a uniformly random function as far as pairwise statistics go, so that the size of $S$ can be read off from the number of hash constraints needed to reduce it to a single point. The isolation lemma uses this to find a unique solution; Stockmeyer's theorem uses it to estimate the number of solutions; Sipser's original $\BPP$ proof uses it to distinguish a large set from a small one.

In every case the deterministic step that makes the method work is the $\NP$ oracle, or the existential quantifier in Sipser's proof, which checks the fibre. The role of randomness is confined to the choice of hash function, and the choice can be made from a pairwise independent family with $O(m)$ bit descriptions, such as the Toeplitz maps over $\mathrm{GF}(2)$, so the randomness is $O(m)$ bits. This is the smallest amount of randomness in any of the counting methods we have described, and it is of interest that it cannot, as far as is known, be reduced further: a deterministic choice of hash functions that worked for every $S$ recognised by a polynomial size circuit would be a hitting set\index{hitting set} for the family of sets of the form $\{ h : |S \cap h^{-1}(0)| \approx |S| 2^{-k} \}$, and explicit such hitting sets are not known (for arbitrary $S$ they cannot exist).

\section{Karp and Luby: counting DNF solutions}

Stockmeyer's theorem needs an $\NP$ oracle because the sets it counts are arbitrary. For some sets the oracle can be dispensed with and the counting done by a randomized polynomial time approximation scheme outright. The first and most instructive example is the number of satisfying assignments of a formula in disjunctive normal form, due to Karp and Luby in 1983.

Let $\varphi = C_1 \vee \dots \vee C_k$ be a DNF formula on $m$ variables, each $C_i$ a conjunction of literals. The number of assignments satisfying $C_i$ alone is $2^{m - |C_i|}$, known exactly. The number satisfying $\varphi$ is $|\bigcup_i S_i|$ where $S_i$ is the set satisfying $C_i$, and inclusion exclusion is exponential. Karp and Luby's method is to sample from the disjoint union $\{ (i, r) : r \in S_i \}$, whose size $\sum_i |S_i|$ is known, by choosing $i$ with probability proportional to $|S_i|$ and then $r$ uniformly in $S_i$, and to accept $(i, r)$ if $i$ is the smallest index with $r \in S_i$. The acceptance probability is $|\bigcup S_i| / \sum |S_i|$, which is at least $1/k$ since each $r$ is counted at most $k$ times in the sum, and so is estimated with relative error $\epsilon$ from $O(k / \epsilon^2)$ samples. Karp, Luby and Madras improved the constants in 1989.

This is the embedding principle of Chapter~\ref{ch:rare} in its purest form. The set to count is embedded in a larger set whose size is known and from which one can sample, and the ratio is not small. The queens embed in the rooks, and the ratio is small. The difference between DNF counting, which has an FPRAS, and queens completion counting, which is $\sharpP$ complete, is exactly that the embedding ratio is bounded below by $1/\mathrm{poly}$ in the first case and is $e^{-\Omega(N)}$ in the second. In the vertical coordinate: for DNF the staircase has its ground step at a width of at least $1/k$, and one draw in $k$ lands on it; for the queens on the empty board the ground step has width $e^{-0.942 N}$ and the ladder is needed.

\section{The complexity of the ratio}

The general principle, which one may extract from Karp and Luby, Jerrum, Valiant and Vazirani, and the vertical analysis together, is this. A self reducible counting problem has an FPRAS\index{FPRAS} if and only if its staircase can be climbed in polynomially many rungs each of which can be estimated by polynomially many samples from a polynomially samplable distribution. The three polynomials correspond to the three sources of difficulty: the length of the ladder, governed by $\log(1/\rho)$ and never the obstruction; the variance at each rung, governed by the ratio and controlled by well balanced schedules; and the cost of a sample, governed by the mixing time. DNF counting has all three polynomial. The permanent, the volume and the ferromagnetic Ising model have all three polynomial, with the third requiring the deep mixing results. The queens completion count has the first two polynomial (the second under the rung ratio hypothesis of Chapter~\ref{ch:prp}) and the third unknown, and the Gent Jefferson Nightingale theorem, with the corollary of Chapter~\ref{ch:prp}, says that if the third is polynomial then $\NP = \RP$.

\section{Toda's theorem and the parity class}

For completeness we sketch why exact counting sits above the hierarchy. The class $\oplus\Pc$ consists of languages $\{ x : f(x) \text{ is odd} \}$ for $f \in \sharpP$. The Valiant Vazirani lemma shows that $\NP \subseteq \RP^{\oplus\Pc}$, since isolating a unique solution and then asking whether the number of solutions is odd decides satisfiability with one sided error. Toda\index{Toda's theorem}'s first step iterates this to show $\PH \subseteq \BPP^{\oplus\Pc}$: each quantifier alternation can be absorbed into the parity oracle at the cost of more randomness. Toda's second step shows that $\BPP^{\oplus\Pc} \subseteq \Pc^{\sharpP}$ by a polynomial identity that converts a parity computation, amplified to have the right parity with high probability, into a count that is divisible by a large power of two in one case and not in the other. The counting oracle then reads off the answer from the count modulo that power.

The theorem has a reading in our terms. The parity of the number of solutions is a property of the staircase, the parity of the width of the ground step, and the hierarchy reduces to it because random hashing can make the ground step have width one with nonnegligible probability, at which point parity is membership. An exact count of the width then determines parity and hence the hierarchy. Approximate counting determines the width to within a factor $(1 \pm \epsilon)$ and so says nothing about parity, which is why it is in the hierarchy while exact counting is not, unless the hierarchy collapses.

\section{Where the queens sit}

We can now locate the queens completion count exactly. Exactly, it is $\sharpP$ complete, and by Toda at least as hard as the whole hierarchy. Approximately, it is in $\BPP^{\NP}$ by Stockmeyer and in $\Sigma_3^p$. Approximately without an oracle, it has an FPRAS if and only if $\NP = \RP$, by Gent, Jefferson and Nightingale\index{Gent, Jefferson and Nightingale} together with the FPRAS to decision direction, which follows from self reducibility. Deterministically and approximately, it has a deterministic polynomial time approximation if and only if $\NP = \Pc$, by the same reasoning without randomness. The vertical walk is an algorithm for the approximate problem that is accurate in polynomial time if a certain Markov chain mixes rapidly and the rung ratios are at least inverse polynomial, and a polynomial level set sampler of some kind exists if and only if $\NP = \RP$.
